\section*{Chapter \ref{ch:b3:astrophysics} --- Astrophysics}

\begin{solution}{exo:b3:astrophysics:1}
(a) $P_{\text{c}} \sim GM^2/R^4 \approx \qty{1.1e15}{Pa}$ --- ten
billion atmospheres. (b) $\bar\rho \approx \qty{1400}{kg/m^3}$:
$T \sim P m_{\text{p}}/\rho k_{\text{B}} \approx \qty{e8}{K}$ ---
crude, an order above the true \qty{1.6e7}{K}, but the scale is
right. (c) Just at the tunnelling threshold. (d) A protostar
contracts and heats \emph{until} the core ignites, then stops
contracting: any gas ball massive enough must arrive at ignition
--- the star tunes itself to the nuclear thermostat.
\end{solution}

\begin{solution}{exo:b3:astrophysics:2}
(a) $L = 4\pi d^2\times1360 \approx \qty{3.8e26}{W}$. (b) $T =
(L/4\pi R^2\sigma)^{1/4} \approx \qty{5800}{K}$. (c) Wien:
$\lambda_{\text{peak}} \approx \qty{500}{nm}$ --- the eye evolved
onto the Sun's peak. (d) Inverse-square flux (Year 1);
Stefan--Boltzmann and Wien: \cref{ch:b3:photons-phonons}'s
blackbody, derived from photon statistics.
\end{solution}

\begin{solution}{exo:b3:astrophysics:3}
(a) Sirius B: far left, far down; Betelgeuse: far right, far up.
(b) $R/R_\odot = \sqrt{L/L_\odot}\,(T_\odot/T)^2$: Sirius B
$\approx 0.009\,R_\odot \approx$ an Earth radius; Betelgeuse
$\approx 860\,R_\odot \approx \qty{4}{au}$. (c) Mercury, Venus,
Earth and Mars would orbit \emph{inside} it. (d) Earth-sized yet
solar-massed: only electron \omterm{thm:b3:quantum-statistics:fermi}{degeneracy pressure}
(\cref{ch:b3:quantum-statistics}) --- Sirius B is the classic
\omterm{ex:b3:quantum-statistics:dwarf}{white dwarf}.
\end{solution}

\begin{solution}{exo:b3:astrophysics:4}
(a) $10\,\text{Gyr}\times10^{-2.5} \approx \qty{30}{Myr}$. (b)
$0.5^{-2.5} \approx 5.7$: about 57 billion years --- every red
dwarf ever born is still burning. (c) Luminosity grows as
$M^{3.5}$ but the tank only as $M$: the rich burn their fortune
faster than they bank it. (d) Massive stars die within a few
million years of birth --- before drifting anywhere: supernovae
detonate inside the star-forming spiral arms that made them.
\end{solution}

\begin{solution}{exo:b3:astrophysics:5}
(a) $E = -K$: radiating makes $E$ more negative, so $K$ ---
i.e.\ the temperature --- \emph{rises}: gravity's thermodynamic
signature. (b) $GM^2/R \approx \qty{3.8e41}{J}$. (c)
$E/L_\odot \approx \qty{e15}{s} \approx \qty{30}{Myr}$. (d)
Kelvin's Sun could not be older than tens of megayears, while
Darwin's worms and Lyell's rocks demanded billions: the
resolution was a fuel nobody knew --- \cref{ch:b3:nuclear-physics},
seventy years early.
\end{solution}

\begin{solution}{exo:b3:astrophysics:6}
(a) $\bar\rho \approx \qty{1.8e9}{kg/m^3}$: a nine-tonne
teaspoon. (b) $g = GM/R^2 \approx \qty{3.3e6}{m/s^2}$ --- three
hundred thousand Earth gravities. (c) $n \propto M/R^3$:
degeneracy supplies $P \propto M^{5/3}/R^5$, gravity demands
$P \propto M^2/R^4$; equating, $R \propto M^{-1/3}$. (d) More
weight needs more Fermi pressure, which only compression
provides --- so mass \emph{shrinks} the star; the spiral has a
wall: once electrons go relativistic the pressure law softens to
$n^{4/3}$, the balance fails at any radius, and
$M_{\text{Ch}} = 1.4\,M_\odot$ is the cliff edge.
\end{solution}

\begin{solution}{exo:b3:astrophysics:7}
(a) $V = M/\rho_{\text{nuc}} \approx \qty{1.2e13}{m^3}$:
$R \approx \qty{14}{km}$ --- a nucleus with a skyline. (b)
$\omega \propto 1/R^2$: spin-up by $(10^8/1.4\times10^4)^2
\approx 5\times10^{7}$, so 25 days shrinks to
${\sim}\qty{40}{ms}$: pulsar territory. (c) $g \approx
\qty{e12}{m/s^2}$: a mountain here would be millimetres there.
(d) $B \approx 10^{-2}\times5\times10^{7} \approx
\qty{5e5}{T}$ --- and the magnetic axis, misaligned with the
spin axis, funnels radio emission along its poles: the beam
sweeps like a lighthouse, and we call the flashes a pulsar.
\end{solution}

\begin{solution}{exo:b3:astrophysics:8}
(a) $\tfrac12mc^2 = GMm/R$ gives $R = 2GM/c^2$. (b) Sun:
\qty{3.0}{km}; Earth: \qty{8.9}{mm} --- compress either inside
that and light cannot leave. (c) $R_{\text{s}} \approx
\qty{1.2e10}{m}$: a fifth of Mercury's orbit --- and stars are
observed whipping around it. (d) $\bar\rho \propto
M/R_{\text{s}}^3 \propto 1/M^2$: a $10^8\,M_\odot$ hole is less
dense than water. Nothing need be ``crushed'' at the horizon ---
it is a point of no return, not a surface of rock.
\end{solution}

\begin{solution}{exo:b3:astrophysics:9}
(a) $H_0 = 7\times10^4/3.09\times10^{22} \approx
\qty{2.3e-18}{s^{-1}}$. (b) $1/H_0 \approx \qty{4.4e17}{s}
\approx 14$ billion years. (c) $v = \qty{1.4e4}{km/s}$;
$z \approx v/c \approx 0.047$: wavelengths arrive stretched
5\,\%. (d) The rate has not been constant --- decelerated by
matter early on, accelerated by dark energy since: that
$1/H_0$ still lands within a gigayear of the true 13.8 is a
minor cosmic coincidence.
\end{solution}

\begin{solution}{exo:b3:astrophysics:10}
(a) $\lambda_{\text{peak}} = b/T \approx \qty{1.1}{mm}$:
microwaves, as advertised. (b) Stretched by $3000/2.725 \approx
1100$ --- the universe is eleven hundred times larger in every
direction since. Recombination waits at \qty{3000}{K}, far
below $\qty{13.6}{eV}/k_{\text{B}}$, because Saha's balance
(\cref{prop:b3:grand-canonical:saha}) is tilted by two billion
photons per \omterm{def:b3:particle-physics:zoo}{baryon}: the ionising tail keeps hydrogen broken up
long after naive energetics says otherwise. (c) About
$1.6\times10^{9}$ photons per \omterm{def:b3:particle-physics:zoo}{baryon} --- the universe is, by
count, almost entirely light. (d) The CMB is not a place but
the whole sky: every line of sight ends on the last-scattering
fog, so every antenna, aimed anywhere, receives it.
\end{solution}

\begin{solution}{exo:b3:astrophysics:11}
(a) $n/p = \eu^{-1.29/0.8} \approx 0.20$. (b) With $n/p = 1/7$:
14 protons per 2 neutrons; the 2 neutrons grab 2 protons into
one $^4$He (mass 4) leaving 12 lone protons: $4/16 = 25\,\%$ by
mass. (c) No star has had time to make a quarter of everything
helium; finding that floor in the oldest, most metal-poor gas
is the hot early universe caught red-handed --- a Boltzmann
factor written across the sky. (d) There is no stable nucleus
at $A = 5$ or $8$, and the density and temperature were
falling within minutes: the bridge to carbon (three heliums at
once) only opens in the dense, patient cores of \omterm{rem:b3:astrophysics:evolution}{red giants}.
\end{solution}

\begin{solution}{exo:b3:astrophysics:12}
(a) $mv^2/r = GM(r)m/r^2$. (b) $M = v^2r/G \approx
\qty{1.1e42}{kg} \approx 5\times10^{11}\,M_\odot$. (c) Roughly
a tenth: ninety percent of the galaxy neither shines nor
absorbs. (d) Flat $v$ means $M(r) \propto r$ --- the unseen
mass keeps accumulating far beyond the visible disc, an
extended dark halo (density $\sim1/r^2$) in which the luminous
galaxy is only the lit-up core.
\end{solution}

\begin{solution}{pb:b3:astrophysics:1}
\textbf{1.} $t \sim 10\,\text{Gyr}\times20^{-2.5} \approx
\qty{6}{Myr}$: the Sun will live two thousand of its lifetimes.
\textbf{2.} Heavier nuclei carry more charge: higher Coulomb
barriers need hotter cores for the tunnelling odds
(\cref{exo:b3:nuclear-physics:8}) to keep pace.
\textbf{3.} Iron sits atop the $B/A$ curve: fusing it
\emph{costs} energy --- the furnace's fuel ladder ends on ash.
\textbf{4.} The \omterm{prop:b3:astrophysics:compact}{Chandrasekhar mass}: the most electron
\omterm{thm:b3:quantum-statistics:fermi}{degeneracy pressure} can carry; the core is a \omterm{ex:b3:quantum-statistics:dwarf}{white dwarf} grown
inside a living star, and at $1.4\,M_\odot$ its support fails.
\textbf{5.} Hydrogen: megayears; helium: hundreds of millennia;
silicon: a day --- each stage pays less per kilogram while the
star, ever hotter, spends (mostly into neutrinos) ever faster:
the clock runs logarithmically mad.
\textbf{6.} $t \sim 1/\sqrt{G\rho} =
1/\sqrt{6.7\times10^{-11}\times10^{12}} \approx \qty{0.1}{s}$:
the core drops out from under the star in a tenth of a second.
\textbf{7.} $E \sim GM^2/R \approx \qty{4e46}{J}$.
\textbf{8.} The photon show is $10^9\times3.8\times10^{26}
\times2.6\times10^{6} \approx \qty{e42}{J}$ --- one part in ten
thousand. The rest leaves as neutrinos: collapsing nuclear
matter is opaque to everything else, but
(\cref{ch:b3:particle-physics}) almost transparent to the weak
force's own particles, which drain the energy in seconds.
\textbf{9.} $E/4\pi d^2 = 4\times10^{46}/2.8\times10^{43}
\approx \qty{1.4e3}{J/m^2}$.
\textbf{10.} At $\qty{15}{MeV} = \qty{2.4e-12}{J}$ each:
${\sim}6\times10^{14}$ neutrinos per square metre --- through
every square metre of Earth.
\textbf{11.} $6\times10^{14}\times10^{-46}\times1.5\times
10^{32} \approx 9$: the detector logged eleven --- an
order-of-magnitude bullseye for a star that died 160\,000 years
ago.
\textbf{12.} The neutrinos leave the core at collapse; the
light is born only when the shock bursts through the star's
envelope hours later. That both arrived the same night after
160 millennia in flight bounds the neutrino to light speed
within parts in $10^{9}$ --- and its mass to almost nothing.
\textbf{13.} The energy budget ($\sim10^{46}\,\unit{J}$, 99\,\%
in neutrinos) and the ${\sim}10$-second burst duration ---
both matched theory's collapsing core, promoting the
supernova mechanism from blackboard to observation.
\textbf{14.} $\bar\rho = 2.8\times10^{30}/(\tfrac43\pi
(1.2\times10^4)^3) \approx \qty{4e17}{kg/m^3}$: nuclear
density --- \cref{def:b3:nuclear-physics:nuclide}'s teaspoon,
now the size of a city.
\textbf{15.} Spin-up $(10^8/1.2\times10^4)^2 \approx
7\times10^{7}$: 30 days $\to$ ${\sim}\qty{40}{ms}$.
\textbf{16.} $B \approx 10^{-2}\times7\times10^{7} \approx
\qty{7e5}{T}$ --- a hundred million tesla-metres of flux,
conserved into a postage stamp.
\textbf{17.} The magnetic axis, tilted off the spin axis,
beams radio along its poles; the beam sweeps the sky like a
lighthouse. The Crab's remnant --- the 1054 guest star --- still
turns thirty times a second, a millennium of spin-down barely
begun.
\textbf{18.} $R_{\text{s}} = 3\times\qty{3}{km} \approx
\qty{9}{km}$ --- of the order of the star itself: with no
pressure left to argue, collapse continues through the horizon
and the remnant is a \omterm{prop:b3:astrophysics:compact}{black hole}.
\textbf{19.} Arrays of such clocks are gravitational-wave
detectors of galactic aperture (and relativity's best
laboratories: binary-pulsar orbits decay exactly as
gravitational radiation predicts).
\textbf{20.} Its oxygen was fused in a massive star's shell
and thrown out by a supernova; the iron was forged in the
explosive end itself; both drifted the interstellar clouds,
joined the solar nebula, the Earth, the sea, the food chain
--- and the reader. Big-bang hydrogen, though: that has been
yours all along.
\textbf{21.} Same detonation mass $\to$ same intrinsic
brightness $\to$ apparent brightness reads distance: a candle
of known wattage visible across billions of light-years ---
the ruler that draws the Hubble diagram's far end.
\textbf{22.} The expansion is \emph{accelerating}: some
``dark energy'' pushes the universe apart --- 2011's Nobel,
and physics' largest open invoice.
\textbf{23.} $\qty{50}{kpc} \approx 163000$ light-years: the
light left while \emph{Homo sapiens} was first leaving Africa.
\textbf{24.} $6\times10^{14}\times0.03 \approx 2\times10^{13}$
through every human alive --- and, at
\cref{exo:b3:particle-physics:10}'s odds, not one person in a
million interacted with even one: the flash that outshone a
galaxy passed through humanity unfelt.
\textbf{25.} A star balanced ten million years on tunnelling
protons; a tenth-of-a-second collapse paid out $10^{46}$
joules in neutrinos; Pauli's principle holds the corpse at
nuclear density; and the scattered debris --- oxygen, iron,
reader --- rides Hubble's expansion still, in a universe whose
baby picture is a Planck curve. Physics, complete for now:
volume closed.
\end{solution}
